% !TeX root = ../../tesis.tex
%%%%%%%%%%%%%%%%%%%%%%%%%%%%%%%%%%%%%%%%%%%%%%%%%%%%%%%%%%%%%%%%%%%
% SECCIÓN 5.5 — ANÁLISIS DE RED: FASE 3
%%%%%%%%%%%%%%%%%%%%%%%%%%%%%%%%%%%%%%%%%%%%%%%%%%%%%%%%%%%%%%%%%%%
\section{Análisis de Red: Fase~3}
\label{sec:fase3-analisis-red}

Las Fases~1 y 2 describen la estructura y los pesos de la red; la
Fase~3 responde una pregunta distinta: ¿cuánto tarda realmente un
viaje en esta red, qué nodos sostienen ese tiempo y qué ocurre cuando
alguno falla? Los tres indicadores de esta fase —tiempo de viaje
promedio, centralidad de intermediación y vulnerabilidad estructural—
operan sobre el componente gigante del grafo (10,561 nodos, 95\,\% de
la red) con el algoritmo de \dijkstra{} ponderado.

% !TeX root = ../../tesis.tex
%%%%%%%%%%%%%%%%%%%%%%%%%%%%%%%%%%%%%%%%%%%%%%%%%%%%%%%%%%%%%%%%%%%
% SECCIÓN 5.5.1 — TIEMPO DE VIAJE PROMEDIO (T)
%%%%%%%%%%%%%%%%%%%%%%%%%%%%%%%%%%%%%%%%%%%%%%%%%%%%%%%%%%%%%%%%%%%
\subsection{Tiempo de Viaje Promedio (\texorpdfstring{$T$}{T})}
\label{subsec:tiempo-viaje-promedio}

Ejecutado el algoritmo de \dijkstra{} ponderado sobre el componente
gigante, el tiempo de viaje promedio de la red resulta en
$T = 108.92$ minutos (corrida: 2026-09-02, 10,561~nodos, tiempo de
cómputo: 251.8~segundos). El valor sintetiza el costo temporal medio
de todos los pares origen-destino posibles bajo las condiciones de
ponderación de Fase~2.

La referencia de campo más cercana proviene de la Encuesta Nacional de
Ocupación y Empleo (ENOE), que estima en aproximadamente 85~minutos el
tiempo promedio de traslado en la \zmvm{}. Frente a los 108.92~minutos
del modelo, la brecha de $\sim$24~minutos tiene origen metodológico, no
en una sobreestimación del sistema. Mientras la ENOE construye su
promedio a partir de desplazamientos observados —concentrados en los
corredores de alta afluencia que conectan el centro con la periferia
próxima— el modelo asigna el mismo peso a cada par del componente
gigante. Los trayectos entre periferias opuestas, que la demanda
cotidiana apenas realiza, quedan contabilizados en igualdad de
condiciones que los recorridos más frecuentes, elevando el promedio
topológico por encima del promedio ponderado por demanda.

Esa variación queda desglosada por alcaldía en la
Figura~\ref{fig:grafica-t-promedio}.

\begin{figure}[H]
    \centering
    \includegraphics[width=0.92\textwidth]{%
        Figures/Mapas/Tableau_Resources/RedActual/VFTModel/UniGrafica/%
        Grafica_VFTModel_2026_Tiempo_Promedio}
    \caption[Tiempo de viaje promedio por alcaldía --- \zmvm{}]{%
    Tiempo de viaje promedio ($T$) por alcaldía. Corrida sobre el
    componente gigante (10,561~nodos) con \dijkstra{} ponderado.
    \zmvm{}, 2026-09-02.}
    \label{fig:grafica-t-promedio}
    \fuentefigura{Elaboración propia con \textbf{Tableau};
    datos: \texttt{VFTModel} \citep{cabrera2026VFTModel}.}
\end{figure}

Con $T$ calculado, surge una pregunta estructural: ¿qué nodos
concentran mayor responsabilidad sobre ese valor? No todos los nodos
son equivalentes en la topología del grafo; algunos articulan una
fracción desproporcionada de las rutas mínimas y su fallo impacta los
tiempos en toda la red. La centralidad de intermediación $B(v)$
—analizada en la sección siguiente— identifica esos puntos.

% !TeX root = ../../tesis.tex
%%%%%%%%%%%%%%%%%%%%%%%%%%%%%%%%%%%%%%%%%%%%%%%%%%%%%%%%%%%%%%%%%%%
% SECCIÓN 5.5.2 — CENTRALIDAD DE INTERMEDIACIÓN B(v)
%%%%%%%%%%%%%%%%%%%%%%%%%%%%%%%%%%%%%%%%%%%%%%%%%%%%%%%%%%%%%%%%%%%
\subsection{Centralidad de Intermediación (\texorpdfstring{$B(v)$}{B(v)})}
\label{subsec:fase3-centralidad}

No todos los nodos de la red tienen el mismo peso en la conectividad
global. Algunos funcionan como puentes inevitables: la mayoría de las
rutas mínimas entre cualquier par de la red pasan por ellos, y su
ausencia fragmentaría la estructura o alargaría los tiempos de forma
generalizada. La centralidad de intermediación $B(v)$ mide exactamente
esa condición: la proporción de rutas mínimas del grafo que atraviesan
cada nodo.

Dado un grafo $\mathcal{G}$ con nodos $\mathcal{V}$, la centralidad de
intermediación de un nodo $v$ se define como:

\begin{equation*}
    B(v) = \sum_{s \neq v \neq t} \frac{\sigma_{st}(v)}{\sigma_{st}}
    \tag{\ref{eq:betweenness}}
\end{equation*}

\noindent donde:
\begin{itemize}
    \item $\sigma_{st}$: cantidad de rutas mínimas entre los nodos $s$ y $t$.
    \item $\sigma_{st}(v)$: subconjunto de esas rutas que pasan por $v$.
\end{itemize}

El modelo calcula $B(v)$ mediante el algoritmo de Brandes, que reduce la
complejidad computacional a $\mathcal{O}(VE)$ sobre grafos ponderados; la
corrida sobre el componente gigante (10,561~nodos) tomó 593.4~segundos.

Dicho de otro modo: hay estaciones que sostienen la red desde adentro
sin que los usuarios lo noten. No son necesariamente las más
transitadas; son las que aparecen en la ruta cuando no hay otra
opción. Quitarlas no molesta a unos cuantos pasajeros —desorganiza
trayectos que atraviesan cuadrantes enteros de la metrópoli.

El Cuadro~\ref{tab:bv-top10} presenta los diez nodos con mayor $B(v)$
calculados sobre el componente gigante. El resultado más llamativo:
los diez son estaciones del Metro, y Tacubaya encabeza la lista con
un valor de 0.2179.

\begin{table}[H]
    \centering
    \caption[Top 10 nodos por centralidad de intermediación $B(v)$]{%
    Diez nodos con mayor centralidad de intermediación ($B(v)$) en la
    red de transporte de la \zmvm{}. Corrida: 2026-09-02.}
    \label{tab:bv-top10}
    \small
    \begin{tabular}{|r|l|l|l|r|}
        \hline
        \textbf{\#} & \textbf{Estación} & \textbf{Sistema} &
        \textbf{Alcaldía} & \textbf{$B(v)$} \\
        \hline
         1 & Tacubaya        & Metro & Miguel Hidalgo & 0.2179 \\
         2 & Mixcoac         & Metro & Benito Juárez  & 0.1725 \\
         3 & Hidalgo         & Metro & Cuauhtémoc     & 0.1595 \\
         4 & Balderas        & Metro & Cuauhtémoc     & 0.1433 \\
         5 & Lázaro Cárdenas & Metro & Cuauhtémoc     & 0.1376 \\
         6 & Chilpancingo    & Metro & Cuauhtémoc     & 0.1366 \\
         7 & Viaducto        & Metro & Benito Juárez  & 0.1346 \\
         8 & Villa de Cortés & Metro & Benito Juárez  & 0.1341 \\
         9 & Nativitas       & Metro & Benito Juárez  & 0.1299 \\
        10 & Zapata          & Metro & Benito Juárez  & 0.1283 \\
        \hline
    \end{tabular}
    \fuentefigura{Elaboración propia con \texttt{VFTModel}
    \citep{cabrera2026VFTModel}. Corrida: 2026-09-02.}
\end{table}

El Metro monopoliza los diez primeros puestos pese a concentrar
menos del 2\,\% de las aristas de servicio real: su ventaja no es
volumétrica sino posicional ---sus estaciones ocupan los puntos donde
más rutas se cruzan.

La distribución geográfica del top~10 es igualmente estrecha: siete
de los diez nodos quedan en Benito Juárez o Cuauhtémoc, lo que revela
que la intermediación de la red está concentrada en un corredor
reducido del centro-poniente de la cuenca.

La Figura~\ref{fig:grafica-bv-nodos} extiende ese ranking a los
principales nodos de la red.

\begin{figure}[H]
    \centering
    \includegraphics[width=0.92\textwidth]{%
        Figures/Mapas/Tableau_Resources/RedActual/VFTModel/UniGrafica/%
        Grafica_VFTModel_2026_Centralidad_Intermediacion_Nodos}
    \caption[Centralidad de intermediación $B(v)$ --- top nodos de la red]{%
    Ranking de los nodos con mayor centralidad de intermediación
    ($B(v)$) en la red de transporte de la \zmvm{}.
    Corrida: 2026-09-02.}
    \label{fig:grafica-bv-nodos}
    \fuentefigura{Elaboración propia con \textbf{Tableau};
    datos: \texttt{VFTModel} \citep{cabrera2026VFTModel}.}
\end{figure}

La centralidad topológica no siempre coincide con la afluencia real.
Para verificarlo, la Figura~\ref{fig:grafica-top10-estaciones} muestra
las diez estaciones con mayor volumen de usuarios según los registros
de Apímetro ---un contraste directo con el ranking de $B(v)$ del
Cuadro~\ref{tab:bv-top10}.

\begin{figure}[H]
    \centering
    \includegraphics[width=0.92\textwidth]{%
        Figures/Mapas/Tableau_Resources/RedActual/Apimetro/UniGrafica/%
        Grafica_Apimetro_2026_Top_10_estaciones}
    \caption[Top 10 estaciones por afluencia de usuarios --- \zmvm{}]{%
    Diez estaciones con mayor volumen de usuarios registrado por
    \texttt{Apímetro} en la red de transporte de la \zmvm{}.
    Corrida: 2026-04-18.}
    \label{fig:grafica-top10-estaciones}
    \fuentefigura{Elaboración propia con \textbf{Tableau};
    datos: \texttt{Apímetro} \citep{cabrera2025apimetro}.}
\end{figure}

La Figura~\ref{fig:grafica-bv-sistema} agrega esos valores por sistema,
mostrando cuánto peso acumulado concentra cada modo en la intermediación
global de la red.

\begin{figure}[H]
    \centering
    \includegraphics[width=0.92\textwidth]{%
        Figures/Mapas/Tableau_Resources/RedActual/VFTModel/UniGrafica/%
        Grafica_VFTModel_2026_Centralidad_Intermediacion_sistema}
    \caption[Centralidad de intermediación $B(v)$ acumulada por sistema]{%
    Suma acumulada de $B(v)$ agrupada por sistema de transporte.
    El Metro concentra la mayor proporción de intermediación global
    de la red. Corrida: 2026-09-02.}
    \label{fig:grafica-bv-sistema}
    \fuentefigura{Elaboración propia con \textbf{Tableau};
    datos: \texttt{VFTModel} \citep{cabrera2026VFTModel}.}
\end{figure}

Para traducirlo al lenguaje urbano: imaginemos que cada par de colonias
de la \zmvm{} traza su ruta más rápida hacia el destino. Las estaciones
que aparecen con mayor frecuencia en ese ejercicio colectivo no son
necesariamente las más concurridas en términos de afluencia, sino las
más difíciles de sustituir: su ausencia obligaría a miles de trayectos
a buscar desvíos más largos o, en el peor caso, a no encontrarlos.
$B(v)$ mide exactamente esa indispensabilidad estructural.

La Figura~\ref{fig:mapa-betweenness} traslada ese ranking al territorio,
mostrando cómo la intermediación topológica se distribuye espacialmente
sobre la cuenca metropolitana.

\begin{figure}[H]
    \centering
    \includegraphics[width=0.92\textwidth]{%
        Figures/Mapas/Qgis_Resources/RedActual/Mapa_Betweenness}
    \caption[Mapa de centralidad de intermediación $B(v)$ --- \zmvm{}]{%
    Distribución espacial de la centralidad de intermediación ($B(v)$)
    en la red de transporte de la \zmvm{}. El corredor de alta
    intermediación traza una banda en el centro-poniente de la cuenca.
    Corrida: 2026-09-02.}
    \label{fig:mapa-betweenness}
    \fuentefigura{Elaboración propia con \textbf{QGIS};
    datos: \texttt{VFTModel} \citep{cabrera2026VFTModel}.}
\end{figure}

La concentración que revelan los datos tabulares adquiere forma geográfica
en el mapa: el corredor de alta intermediación no cubre la metrópoli de
manera uniforme, sino que traza una banda relativamente estrecha sobre el
centro-poniente, dejando amplias zonas periféricas con valores bajos o
nulos. Esas zonas ---las mismas que registraron menor cobertura y mayor
Factor de Desviación en Fase~1--- son también las que dependen de los
nodos del corredor central para cualquier viaje de largo aliento.

La Figura~\ref{fig:dashboard-bv} consolida la visión integral del indicador
sobre el conjunto de la red.

\begin{figure}[H]
    \centering
    \includegraphics[width=0.92\textwidth]{%
        Figures/Mapas/Tableau_Resources/RedActual/VFTModel/Dashboards/%
        Dashboard_VFTModel_2026_Centralidad_intermediacion}
    \caption[Dashboard integral de centralidad de intermediación $B(v)$ --- \zmvm{}]{%
    Vista integral del indicador de centralidad de intermediación ($B(v)$)
    sobre la red de transporte de la \zmvm{}.
    Corrida: 2026-09-02.}
    \label{fig:dashboard-bv}
    \fuentefigura{Elaboración propia con \textbf{Tableau};
    datos: \texttt{VFTModel} \citep{cabrera2026VFTModel}.}
\end{figure}

Los tres indicadores de Fase~3 ---tiempo de viaje, intermediación y
vulnerabilidad preliminar--- apuntan en la misma dirección: la red
concentra su capacidad de articulación en un corredor reducido del
centro-poniente, mientras la periferia sur y oriente queda
estructuralmente subordinada a él. La sección siguiente examina qué
ocurre cuando alguno de esos nodos críticos desaparece de la red.

% !TeX root = ../../tesis.tex
%%%%%%%%%%%%%%%%%%%%%%%%%%%%%%%%%%%%%%%%%%%%%%%%%%%%%%%%%%%%%%%%%%%
% SECCIÓN 5.5.3 — VULNERABILIDAD ESTRUCTURAL ΔE
%%%%%%%%%%%%%%%%%%%%%%%%%%%%%%%%%%%%%%%%%%%%%%%%%%%%%%%%%%%%%%%%%%%
\subsection{Vulnerabilidad Estructural (\texorpdfstring{$\Delta E$}{Delta E})}
\label{subsec:vulnerabilidad-estructural}

Los resultados de esta subsección tienen carácter preliminar. La
simulación de vulnerabilidad estructural se realizó sobre los nodos
del grafo construido por \texttt{VFTModel}, pero no fue ejecutada
sobre la capa de red integrada de \texttt{Transport-gis-mjg}; en
consecuencia, los valores reportados no han sido verificados contra
la representación completa del sistema. Se presentan como primer
acercamiento al indicador $\Delta E$, con la intención de enmarcar
el método y señalar las líneas de trabajo futuro.

A partir del valor de referencia $T = 108.92$~min
(§\ref{subsec:tiempo-viaje-promedio}), el experimento calcula la
caída relativa de eficiencia $\Delta E$ cuando se retira cada uno de
los cinco nodos de mayor $B(v)$ (§\ref{subsec:robustez-geometrica}).
La eficiencia base es $E_0 = 1/108.92$; tras remover el nodo $v$,
se recalcula el tiempo promedio $T_f$ y se obtiene
$\Delta E = (E_0 - E_f)/E_0$.

El Cuadro~\ref{tab:delta-e-top5} registra los resultados del
experimento sobre los cinco nodos seleccionados.

\begin{table}[H]
    \centering
    \caption[Impacto de remoción de nodos críticos --- $\Delta E$ exploratorio]{%
    Caída relativa de eficiencia $\Delta E$ al remover cada nodo del
    componente gigante. $E_0 = 1/108.92$; $E_f = 1/T_f$. Corrida
    exploratoria: 2026-09-02.}
    \label{tab:delta-e-top5}
    \small
    \begin{tabular}{|r|l|r|r|r|}
        \hline
        \textbf{\#} & \textbf{Nodo removido} &
        \textbf{$B(v)$} & \textbf{$T_f$ (min)} & \textbf{$\Delta E$} \\
        \hline
        1 & Tacubaya        & 0.2179 & 111.80 & 2.58\,\% \\
        2 & Mixcoac         & 0.1725 & 110.19 & 1.15\,\% \\
        3 & Hidalgo         & 0.1595 & 109.90 & 0.89\,\% \\
        4 & Balderas        & 0.1433 & 109.57 & 0.59\,\% \\
        5 & Lázaro Cárdenas & 0.1376 & 110.13 & 1.10\,\% \\
        \hline
    \end{tabular}
    \fuentefigura{Elaboración propia con \texttt{VFTModel}
    \citep{cabrera2026VFTModel}. Corrida exploratoria: 2026-09-02.}
\end{table}

Tacubaya encabeza tanto el ranking de $B(v)$ como el de $\Delta E$: su
remoción eleva el tiempo promedio de la red en 2.88~min, traducido en
una caída de eficiencia del 2.58\,\%. El dato es coherente con la
posición estructural del nodo ---cruce de las Líneas 1, 7 y 9 del
Metro--- y su papel de puente entre el corredor poniente y el centro.
Los cuatro nodos restantes producen caídas entre 0.59\,\% y 1.15\,\%,
lo que confirma que ningún nodo individual colapsa la red, pero la
pérdida acumulada de varios podría alcanzar umbrales críticos.

Los tres indicadores de Fase~3 convergen en un mismo diagnóstico
estructural. $T = 108.92$~min revela que el tiempo de viaje promedio
entre cualquier par de la \zmvm{} supera con amplitud los umbrales de
movilidad cotidiana. $B(v)$ identifica el corredor que hace posible
ese tiempo: un conjunto reducido de estaciones del Metro en el eje
centro-poniente concentra la mayor parte de los caminos mínimos de la
red. Y la simulación exploratoria de $\Delta E$ confirma que la
remoción de esos nodos, aunque no colapsa la red, la degrada de
forma mensurable ---con Tacubaya produciendo el mayor impacto
individual registrado.

El patrón es consistente con el hallazgo de Fase~1: las zonas que
registraron menor cobertura y mayor Factor de Desviación son también
las que más dependen del corredor central para cualquier trayecto de
largo aliento. La siguiente sección integra estos indicadores en un
diagnóstico territorial que localiza las zonas de déficit estructural
de la red.

